\section{Economic Equilibrium}\label{app:EE}\setcounter{equation}{0}
The appendix derives results for the economic equilibrium with CRRA-GHH and log-GHH preferences and where informal workers' endowment is tied to wages in the formal sector absent taxes. 
\subsection{With CRRA-GHH preferences}\label{app:EE:CRRA}
\smalltitle{Household solutions}
With CRRA-GHH preferences, the first order conditions in \eqref{eq:model:formalFOCs} can be written as
\begin{align*}
    h_{t,i} &= \left(\dfrac{\eta_iw_t(1-\tau_t)+\frac{\theta \overline{b}_{t+1} \eta_i}{R_{t+1}}}{X_i}\right)^{\xi} \\
    c_{2,t+1}^i &= \tilde{c}_{1,t}^i\left(\beta R_{t+1}\right)^{\rho}.
\end{align*}
Using the budget constraints and rearranging slightly, we can write individual savings as
\begin{align*}
    s_{t,i} = \dfrac{B_{t+1}}{1+B_{t+1}}\left[w_t\left(1-\tau_t\right)h_{t,i}\eta_i-\dfrac{X_{i}\xi}{1+\xi}\left(h_{t,i}\right)^{\frac{1+\xi}{\xi}}\right]-\dfrac{b_{t+1}^i/R_{t+1}}{1+B_{t+1}},
\end{align*}
where we have used the shorthand notation $B_{t+1} \equiv \left(\beta \right)^{\rho}(R_{t+1})^{\rho-1}$.


When informal households save using informal technology, similar conditions apply:\footnote{Note that given our assumption on universal pensions, informal workers do not perceive a link between their work effort and the benefits they might receive.}
\begin{align}
    h_{t,0} &= \left(\dfrac{\eta_{0}w_t^0}{X_{0}}\right)^{\xi}\label{here1} \\
    c_{2,t+1}^0 &= \tilde{c}_{1,t}^0\left(\beta R_{t+1}^0\right)^{\rho}. 
\end{align}
As for formal household we can derive the following for informal savings:
\begin{align*}
    s_{t,0} = \dfrac{B_{t+1}^0}{1+B_{t+1}^0}\left[w_t^0h_{t,0}\eta_{0}-\dfrac{X_{0}\xi}{1+\xi}\left(h_{t,0}\right)^{\frac{1+\xi}{\xi}}\right]-\dfrac{b_{t+1}^0/R_{t+1}^0}{1+B_{t+1}^0},
\end{align*}
where we have used $B_{t+1}^0 \equiv \left(\beta \right)^{\rho}(R_{t+1}^0)^{\rho-1}$. When informal households are hand-to-mouth and receive an endowment when old, the solution for labor supply is given by (\ref{here1}).
%\begin{align*}
 %   h_{t,0} = \left(\dfrac{\eta_{0}w_t^0}{X_{0}}\right)^{\xi} %\\
  %  h_{t+1,0}^r &= \left(\dfrac{\chi \eta_{0}w_{t+1}^0}{X_{0}}\right)^{\xi}
%\end{align*}




\smalltitle{Savings and labor supply in equilibrium.}
Next, in economic equilibrium, we impose conditions for factor prices in equations \eqref{eq:factorPrices} and the specification of pension benefits including the balanced budget condition in equations \eqref{eq:model:governmentBudget}. We can then rewrite type $i$'s labor supply decision as
\begin{align*}
    h_{t,i} = \left(\dfrac{\eta_i}{X_i}\right)^{\xi}\left(\dfrac{1}{h_t}\right)^{\xi}\left((1-\alpha)(1-\tau_t)\left(\dfrac{s_{t-1}}{\nu_t}\right)^{\alpha}h_t^{1-\alpha}+\dfrac{1-\alpha}{\alpha}\dfrac{\theta}{1+\gamma_0 \epsilon} s_t\tau_{t+1}\right)^{\xi}.
\end{align*}


The only type-specific term is $(\eta_i/X_i)^{\xi}$; this shows that the ratio $h_{t,i}/h_t$ is only a function of parameters, namely that
\begin{align}\label{eq:hi_h}
    \dfrac{h_{t,i}}{h_t} = \dfrac{(\eta_i/X_i)^{\xi}}{\Gamma_h}, \qquad \Gamma_{h} \equiv \sum_{i>0}\frac{\gamma_{i}\eta_{i}^{1+\xi}}{X_{i}^{\xi}}.
\end{align}

Next, we can use the individual labor decision and wage rates to write 
\begin{align*}
    (1-\alpha)(1-\tau_t)h_t^{1-\alpha}\left(\dfrac{s_{t-1}}{\nu_t}\right)^{\alpha} = \frac{h_t^{\frac{1+\xi}{\xi}}}{\Gamma_{h}^{1/\xi}} - \dfrac{1-\alpha}{\alpha}\dfrac{\theta}{1+\gamma_0\epsilon} s_t\tau_{t+1},
\end{align*}
Using this and the discounted pension benefits in equilibrium, we have that
\begin{align}
    s_{t,i} =& \dfrac{1}{1+B_{t+1}}\left[B_{t+1}\frac{\eta_{i}^{1+\xi}}{X_{i}^{\xi}}\left(\frac{h_t}{\Gamma_{h}}\right)^{\frac{1+\xi}{\xi}}\frac{1}{1+\xi}-\frac{1-\alpha}{\alpha}\frac{1-\theta}{1+\gamma_0\epsilon}s_t\tau_{t+1}\right] \nonumber \\
    -&\frac{\eta_{i}^{1+\xi}}{X_{i}^{\xi}\Gamma_{h}}\frac{1-\alpha}{\alpha}\frac{\theta}{1+\gamma_0\epsilon}s_t\tau_{t+1} \label{here2}
\end{align}
Using the definition of the aggregate state \eqref{eq:model:aggStates} and rearranging, we then have:
\begin{align}\label{eq:s}
    s_t = \left(\frac{h_t}{\Gamma_{h}}\right)^{\frac{1+\xi}{\xi}} \Gamma_{s,t}, \qquad \Gamma_{s,t} \equiv \frac{1}{1+\xi}\frac{\Gamma_h B_{t+1} }{1+B_{t+1}+\frac{1-\alpha}{\alpha}\frac{\tau_{t+1}}{1+\gamma_0\epsilon}\left(1+\theta B_{t+1}\right)}
\end{align}
Furthermore, note that we can combine equations for $s_t, h_t$ to yield a closed form expressions \textit{given} values of $B_{t+1}$. From the definition of $h_t$ we use that $s_t$ is proportional to $h_t^{(1+\xi)/\xi}$ to derive:
\begin{align}\label{eq:h}
    h_t = \Theta_{h,t} \left(\frac{s_{t-1}}{\nu_t}\right)^{\frac{\alpha\xi}{1+\alpha\xi}}, \qquad \Theta_{h,t} \equiv \left(\Gamma_{h}\right)^{\frac{1+\xi}{1+\alpha\xi}}\left(\frac{(1-\alpha)(1-\tau_t)}{\Gamma_{h}-\frac{1-\alpha}{\alpha}\frac{\theta\tau_{t+1}}{1+\gamma_0\epsilon}\Gamma_{s,t}}\right)^{\frac{\xi}{1+\alpha\xi}}
\end{align}
Using this, an alternative expression for $s_t$ is:
\begin{align}
    s_t &= \Theta_{s,t}\left(\frac{s_{t-1}}{\nu_t}\right)^{\frac{\alpha(1+\xi)}{1+\alpha\xi}}, \qquad \Theta_{s,t} \equiv \left(\frac{\Theta_{h,t}}{\Gamma_{h}}\right)^{\frac{1+\xi}{\xi}}\Gamma_{s,t}. 
\end{align}
Using the relation between $s_t$ and $h_t^{\frac{1+\xi}{\xi}}$ in (\ref{eq:h}), the ratio of individual (\ref{here2}) to aggregate savings satisfies:
\begin{align}\label{eq:si_s}
    \frac{s_{t,i}}{s_t} = \frac{\eta_i^{1+\xi}}{X_i^{\xi}} \frac{1}{\Gamma_h} + \frac{1-\alpha}{\alpha} \frac{\tau_{t+1}}{1+\gamma_0\epsilon}\frac{1-\theta}{1+B_{t+1}}\left(\frac{\eta_i^{1+\xi}}{X_i^{\xi}\Gamma_h}-1\right).
\end{align}


\bigskip 
For informal households, we assume that informal wages are proportional to formal wages absent taxes and scaled appropriately by informal productivity $\eta_0$. Furthermore, we assume that the return to informal savings is proportional to the formal return on savings $R_t^0 = r_0 R_t$.


Specifically, we assume that informal wages are given by:
\begin{align*}
    w_t^0 = \left(\frac{1-\alpha}{\Gamma_{h}^{\alpha}}\right)^{\frac{1}{1+\alpha\xi}} \left(\frac{s_{t-1}}{\nu_t}\right)^{\frac{\alpha}{1+\alpha\xi}}
\end{align*}
Informal labor supply then follows directly from the derivations above for formal workers; in the case with informal savings, the ratio of informal to aggregate savings is similarly given by
\begin{align}\label{eq:s0_s}
    \frac{s_{t,0}}{s_t} =& \frac{B_{t+1}^0}{1+B_{t+1}^0}\frac{\eta_{0}^{1+\xi}}{X_{0}^{\xi}(1+\xi)}\left(\frac{1-\alpha}{\Gamma_{h}^{\alpha}}\right)^{\frac{1+\xi}{1+\alpha\xi}}
    \frac{1}{\Theta_{s,t}}-\frac{1}{1+B_{t+1}^0}\frac{1-\alpha}{r_0 \alpha}\frac{\epsilon\tau_{t+1}}{1+\gamma_0\epsilon }
\end{align}


\smalltitle{Consumption functions.}
It is straightforward to derive relevant equilibrium consumption functions ($c_{1,t}^j, \tilde{c}_{1,t}^j, c_{2,t}^j$) given the results above. Imposing household budget constraints (\ref{eq:model:BC}), factor prices \eqref{eq:factorPrices}, pension benefits and balanced budget condition \eqref{eq:model:governmentBudget}, and labor and savings decisions from above, we have for formal households $i$:
\begin{align}
    c_{1,t}^i &= \frac{\eta_{i}^{1+\xi}}{X_{i}^{\xi}}\frac{s_t}{\Gamma_{s,t}}\left(1-\frac{B_{t+1}}{(1+B_{t+1})(1+\xi)}\right)+\frac{s_t}{1+B_{t+1}}\frac{1-\alpha}{\alpha}\frac{ \tau_{t+1}(1-\theta)}{1+\gamma_0\epsilon} \\
    c_{2,t}^i &= \alpha \nu_t h_t^{1-\alpha}\left(\frac{s_{t-1}}{\nu_t}\right)^{\alpha}\left(\frac{s_{t-1,i}}{s_{t-1}}+\frac{1-\alpha}{\alpha}\frac{\tau_t}{1+\gamma_0\epsilon}\left(1+\theta\left[\frac{\eta_{i}^{1+\xi}}{X_{i}^{\xi}}\frac{1}{\Gamma_{h}}-1\right]\right)\right) \\
    \tilde{c}_{1,t}^i &= \frac{s_t}{1+B_{t+1}}\left(\frac{\eta_{i}^{1+\xi}}{X_{i}^{\xi}}\frac{1}{(1+\xi)\Gamma_{s,t}}+\frac{1-\alpha}{\alpha}\frac{\tau_{t+1}(1-\theta)}{1+\gamma_0\epsilon}\right)
\end{align}


When informal households save using informal technology, equilibrium consumption functions are given by
\begin{align}
    \tilde{c}_{1,t}^0 =& \frac{1}{1+\xi}\frac{\eta_{0}^{1+\xi}}{X_{0}^{\xi}}\left(\frac{1-\alpha}{\Gamma_{h}^{\alpha}}\right)^{\frac{1+\xi}{1+\alpha\xi}}\left(\frac{s_{t-1}}{\nu_t}\right)^{\frac{\alpha(1+\xi)}{1+\alpha\xi}}+\frac{1}{1+B_{t+1}^0}\frac{1-\alpha}{r_0 \alpha}\frac{\epsilon\tau_{t+1}}{1+\gamma_0\epsilon} s_t \\ 
    c_{2,t}^0 =& r_0\alpha \nu_t\left(\frac{s_{t-1}}{\nu_t}\right)^{\alpha}h_t^{1-\alpha}\left(\frac{s_{t-1,0}}{s_{t-1}}+\frac{1-\alpha}{r_0\alpha} \frac{\epsilon\tau_t}{1+\gamma_0\epsilon }\right).
\end{align}


When informal households are hand-to-mouth, relevant equilibrium consumption functions are given by\footnote{Recall informal workers are assumed to receive an endowment proportional to wages of formal workers when they retire. For tractability we define this as if they were making a labor supply decision. The relevant parameter determining the size of the endowment is $\chi$.}
\begin{align}
    \tilde{c}_{1,t}^0 &= \frac{1}{1+\xi}\frac{\eta_{0}^{1+\xi}}{X_{0}^{\xi}}\left(\frac{1-\alpha}{\Gamma_{h}^{\alpha}}\right)^{\frac{1+\xi}{1+\alpha\xi}}\left(\frac{s_{t-1}}{\nu_t}\right)^{\frac{\alpha(1+\xi)}{1+\alpha\xi}} \\
    \tilde{c}_{2,t}^0 &= \frac{1}{1+\xi}\frac{(\chi \eta_{0})^{1+\xi}}{X_{0}^{\xi}}\left(\frac{1-\alpha}{\Gamma_{h}^{\alpha}}\right)^{\frac{1+\xi}{1+\alpha\xi}}\left(\frac{s_{t-1}}{\nu_t}\right)^{\frac{\alpha(1+\xi)}{1+\alpha\xi}}+\dfrac{(1-\alpha)\nu_t\epsilon \tau_t}{1+\gamma_0\epsilon}\left(\frac{s_{t-1}}{\nu_t}\right)^\alpha h_t^{1-\alpha}.
\end{align}



\subsection{With log-GHH preferences}\label{app:EE:LOG}
With log-preferences, the economic equilibrium is identical to the CRRA equations outlined in appendix \ref{app:EE:CRRA}, except for the discount function $B_t =B^0_t= \beta$. All relevant variables $x_t\in \lbrace s_t, h_t, s_t^i, h_t^j, \tilde{c}_{1,t}^j, c_{2,t}^i, \tilde{c}_{2,t}^0, c_{2,t+1}^i, \tilde{c}_{2,t+1}^0\rbrace$ can be written on the form $x_t = \Theta_{x,t}\left(\frac{s_{t-1}}{\nu_t}\right)^{\sigma_x}$ where $\sigma_x$ are parameter functions
\begin{align*}
    \sigma_{h} &= \sigma_{h}^j = \dfrac{\alpha \xi}{1+\alpha \xi} \\ 
    \sigma_{\tilde{c}_1,t}^j &= \sigma_{c_2,t}^i = \sigma_{\tilde{c}_2,t}^0 = \sigma_{s,t} = \sigma_{s,t}^i =\dfrac{\alpha(1+\xi)}{1+\alpha\xi} \\
    \sigma_{c_2,t+1}^i &= \sigma_{\tilde{c}_2,t+1}^0 = \left(\dfrac{\alpha(1+\xi)}{1+\alpha\xi}\right)^2,
\end{align*}
and where functions $\Theta_{x,t}$ only depends on parameters and policy. 

\smalltitle{Proof of proposition \ref{prop:LOG:EE}.} Part \emph{i.} outlines the marginal effect on $s_t, h_t$ from policy and pension characteristics. Given the solution structure, the effect on $s_t, h_t$ can be found by differentiating equations \eqref{eq:s} and \eqref{eq:h} with respect to the relevant variable. For $\tau_t$ the effect is straightforward:
\begin{align*}
    \frac{\partial \ln(h_t)}{\partial \tau_t} &= -\frac{\xi}{1+\alpha\xi}\frac{1}{1-\tau_t} \leq 0\\
    \dfrac{\partial \ln(s_t)}{\partial \tau_t} &= -\frac{1+\xi}{1+\alpha\xi} \frac{1}{1-\tau_t} \leq 0.
\end{align*}
We can also show that $h_t$ is unambiguously decreasing in $\epsilon$:
\begin{align*}
    \frac{\partial \ln(h_t)}{\partial \epsilon} &= \frac{\xi}{1+\alpha\xi}\frac{\frac{1-\alpha}{\alpha}\frac{\theta \tau_{t+1}}{1+\gamma_0\epsilon}\Gamma_{s,t}}{\Gamma_h-\frac{1-\alpha}{\alpha}\frac{\theta\tau_{t+1}}{1+\gamma_0\epsilon}\Gamma_{s,t}}\frac{1}{1+\gamma_0\epsilon}\left((1+\gamma_0\epsilon)\frac{\partial \ln\left(\Gamma_{s,t}\right)}{\partial \epsilon}-\gamma_0\right)\leq 0 \\
    \frac{\partial \ln\left(\Gamma_{s,t}\right)}{\partial \epsilon}&= \frac{\gamma_0}{1+\gamma_0\epsilon}\frac{\frac{1-\alpha}{\alpha}\frac{\tau_{t+1}}{1+\gamma_0\epsilon}(1+\theta\beta)}{1+\beta+\frac{1-\alpha}{\alpha}\frac{\tau_{t+1}}{1+\gamma_0\epsilon}(1+\theta \beta)}.
\end{align*} 
However, the effect on $s_t$ is harder to sign. 
\begin{align*}
 \dfrac{\partial \ln\left(s_t\right)}{\partial \epsilon} = \frac{1+\xi}{\xi}\dfrac{\partial \ln\left(h_t\right)}{\partial \epsilon}+ \frac{\partial \ln\left(\Gamma_{s,t}\right)}{\partial \epsilon} . 
 \end{align*}
Replacing from the expressions above, and neglecting common factors, the sign of the derivative is given by the sign of
\begin{align*}
-\frac{1}{1+\alpha \xi}\frac{\theta \beta(1+\beta)}{1+\beta +\frac{1-\alpha}{\alpha}\frac{\tau_{t+1}}{1+\gamma_0 \epsilon}(1+\theta \beta \frac{\xi}{1+\xi})}+1+\theta \beta.
 \end{align*}
This can be seen to be positive for all $\theta \in[0,1]$. Note that the effect from $\gamma_0$ is equivalent to that of $\epsilon$. We now consider the effect of future taxes on labor supply and savings. 
\begin{align*}
    \dfrac{\partial \ln\left(h_t\right)}{\partial \tau_{t+1}} &= \frac{\xi}{1+\alpha\xi} \frac{\frac{1-\alpha}{\alpha}\frac{\theta \Gamma_{s,t}}{1+\gamma_0\epsilon}}{\Gamma_h-\frac{1-\alpha}{\alpha}\frac{\theta \tau_{t+1}}{1+\gamma_0\epsilon}\Gamma_{s,t}}\left(1+\tau_{t+1}\frac{\partial \ln\left(\Gamma_{s,t}\right)}{\partial \tau_{t+1}}\right) \geq 0 \\
    \tau_{t+1} \dfrac{\partial \ln\left(\Gamma_{s,t}\right)}{\partial \tau_{t+1}} &= - \frac{\frac{1-\alpha}{\alpha}\frac{\tau_{t+1}(1+\theta \beta)}{1+\gamma_0\epsilon}}{1+\beta+\frac{1-\alpha}{\alpha}\frac{\tau_{t+1}(1+\theta \beta)}{1+\gamma_0\epsilon}} \in (-1,0].
\end{align*}
This shows that labor supply is increasing in $\tau_{t+1}$. For savings the effect is:
\begin{align*}
 \dfrac{\partial \ln\left(s_t\right)}{\partial \tau_{t+1}} = \frac{1+\xi}{\xi}\dfrac{\partial \ln\left(h_t\right)}{\partial \tau_{t+1}}+ \frac{\partial \ln\left(\Gamma_{s,t}\right)}{\partial \tau_{t+1}} . 
 \end{align*}
As before the sign of this derivative can be shown to be the same as the sign of the following term:
\begin{align*}
\frac{1}{1+\alpha \xi}\frac{\theta \beta(1+\beta)}{1+\beta +\frac{1-\alpha}{\alpha}\frac{\tau_{t+1}}{1+\gamma_0 \epsilon}(1+\theta \beta \frac{\xi}{1+\xi})}-(1+\theta \beta).
 \end{align*}
This is basically the same term (with different sign) as before. Thus, for all values of $\theta$ and $\tau_{t+1}$ we have that an increase in $\tau_{t+1}$ reduces savings. 



Taking the derivative with respect to $\theta$, we see a similar condition as with the effect of $\tau_{t+1}$, namely:
\begin{align*}
    \dfrac{\partial \ln\left(h_t\right)}{\partial \theta} &= \frac{\xi}{1+\alpha\xi} \frac{\frac{1-\alpha}{\alpha}\frac{\tau_{t+1} \Gamma_{s,t}}{1+\gamma_0\epsilon}}{\Gamma_h-\frac{1-\alpha}{\alpha}\frac{\theta \tau_{t+1}}{1+\gamma_0\epsilon}\Gamma_{s,t}}\left(1+\theta\frac{\partial \ln\left(\Gamma_{s,t}\right)}{\partial \theta}  \right) \geq 0, \\
    \theta \dfrac{\partial \ln\left(\Gamma_{s,t}\right)}{\partial \theta} &= - \frac{\frac{1-\alpha}{\alpha}\frac{\tau_{t+1}\theta \beta}{1+\gamma_0\epsilon}}{1+\beta+\frac{1-\alpha}{\alpha}\frac{\tau_{t+1}(1+\theta \beta)}{1+\gamma_0\epsilon}} \in (-1,0]. 
\end{align*}
This shows that labor supply is increasing in $\theta$. For the savings rate, we have: 
\begin{align*}
\dfrac{\partial \ln\left(s_t\right)}{\partial \theta} = \frac{1+\xi}{\xi}\dfrac{\partial \ln\left(h_t\right)}{\partial \theta}+ \frac{\partial \ln\left(\Gamma_{s,t}\right)}{\partial \theta}.
\end{align*}
Proceeding as before we find that this is unambiguously negative.\footnote{The term whose sign determines this is slightly different
\begin{align*}
\frac{1}{1+\alpha \xi}\frac{1+\beta+\frac{1-\alpha}{\alpha}\frac{\tau_{t+1}}{1+\gamma_0 \epsilon}}{1+\beta+\frac{1-\alpha}{\alpha}\frac{\tau_{t+1}}{1+\gamma_0 \epsilon}(1+\theta \beta \frac{\xi}{1+\xi})}-1.
 \end{align*}
 }

Part \emph{ii.} follows from earlier derivations, specifically equation \eqref{eq:hi_h}. Part \emph{iii.} similarly follows from equation \eqref{eq:si_s} with $B_{t+1} = \beta$ in the log case. Finally, part \emph{iv.} follows straightforwardly, comparing \eqref{eq:si_s} and \eqref{eq:hi_h}: An average type $i$ household has $\eta_i^{1+\xi}/X_i^{\xi} = \Gamma_h$. Thus, for $\tau_{t+1}=0$ or $\theta = 1$, we simply have $s_{t,i} / s_t = \eta_ih_{t,i}/h_t$. The spread in savings rates increase in $\tau_{t+1}$ and decrease in $\theta$ in the sense that these parameter changes increase $s_{t,i}/s_t$ for $s_{t,i}/s_t>1$ and decrease for $s_{t,i}/s_t<1$.  

Part \emph{v.} The effects on $s_{t,0}$ and $s_{t,0}/s_t$ from various parameters come from differentiation of \eqref{eq:s0_s} with $B_{t+1}^0 = \beta$. Note that $s_{t,0}$ only depends on $\tau_t$ though the effect of taxes on $s_t$. Thus, it is straightforward to show that both $s_{t,0}$ and $s_{t,0}/s_t$ are increasing in $\tau_t$.

The effects from $\tau_{t+1}$ and $\epsilon$ are given by:
\begin{align*}
    \frac{\partial s_{t,0}}{\partial \tau_{t+1}} &= -\frac{1}{1+\beta}\frac{1-\alpha}{r_0\alpha} \frac{s_t\epsilon}{1+\gamma_0\epsilon}\left(1+ \tau_{t+1}\frac{\partial \ln(s_t)}{\partial \tau_{t+1}}\right) \leq 0, \\
    \frac{\partial s_{t,0}}{\partial \epsilon} &= -\frac{1}{1+\beta}\frac{1-\alpha}{r_0\alpha} \frac{s_t\tau_{t+1}}{1+\gamma_0\epsilon}\left(\frac{1}{1+\gamma_0\epsilon}+ \epsilon \frac{\partial \ln(s_t)}{\partial \epsilon}\right)  \leq 0.  
\end{align*}
Where the first derivative can be signed since $\tau_{t+1}\frac{\partial \ln(s_t)}{\partial \tau_{t+1}} \geq -1$. It is straightforward to show that an increase (decrease) in $\theta$ ($\gamma_0$) increases $s_t,0$. 

To see the effect of $\tau_{t+1}$ on the ratio $s_{t,0}/s_t$ we have that:
\begin{align*}
    \frac{\partial (s_{t,0}/s_t)}{\partial \tau_{t+1}} &= -\frac{\beta}{1+\beta}\frac{\eta_{0}^{1+\xi}}{X_{0}^{\xi}(1+\xi)}\left(\frac{1-\alpha}{\Gamma_{h}^{\alpha}}\right)^{\frac{1+\xi}{1+\alpha\xi}}
    \frac{1}{\Theta_{s,t}} \frac{\partial \ln\left(\Theta_{s,t}\right)}{\partial \tau_{t+1}}-\frac{1}{1+\beta}\frac{1-\alpha}{r_0\alpha}\frac{\epsilon}{1+\gamma_0\epsilon},
\end{align*}
where the derivative $\partial \ln\left(\Theta_{s,t}\right)/\partial \tau_{t+1} = \partial \ln\left(s_t\right)/\partial \tau_{t+1}$ is defined above. We note that this may be negative or positive. For instance, if $\epsilon \approx 0$ an increase in future taxes increases this ratio. Conversely if $r_0 \approx 0$ then we would expect the opposite effect. 


